\section{Introduction}

Beyond transcriptomic profiling, single-cell technologies increasingly aim
to capture complementary molecular layers from the same cell, including
cell-surface protein states that RNA alone does not resolve. Aptamers offer
one route to this: short, selected oligonucleotide probes that bind surface
targets and, because their identity is itself a nucleic-acid sequence, can
be read out and paired with transcriptome sequencing on the same cell.
Apt-seq first demonstrated that jointly measured aptamer-binding and RNA
profiles resolve cell composition in defined cell-line mixtures
\citep{delley2018combined}. Subsequent work extended this along two
directions: Aptomics scaled aptamer-based profiling to proteins, glycans,
and RNA in clinical specimens, capturing tumor heterogeneity and T-cell
differentiation states \citep{wu2025aptomics}, while SPARK-seq combined
CRISPR perturbation with aptamer and RNA sequencing to identify and
kinetically characterize surface-binding targets \citep{luo2026spark}.
Together, these studies establish that APT measurements are feasible at
single-cell resolution, discriminate defined cell populations, and support
downstream biological characterization. What remains less established is
whether APT features on their own support reliable identification of a
cell's annotated identity -- an evaluation that determines what downstream
model development and clinical application this modality can support. As
single-cell APT data accumulates across patient cohorts, this capability
needs to be evaluated directly: which annotated cell identities can APT
features support in a patient who contributed no data to model
development, and does paired APT improve on a specified RNA-only
reference?

Two comparisons are needed to answer this question. First, low fine-subtype
performance does not establish a lack of APT information: parent-lineage
errors, candidate-set size, annotation ambiguity and the fitted model all
affect reconstruction. True-lineage decoding must therefore be compared
with a training-only prior given the same privileged lineage. Second, an
RNA+APT gain can reflect patient structure or coarse-lineage correspondence
rather than finer cell pairing. Shuffling APT within patient and then within
patient and lineage provides progressively more specific controls.

APT-Bench brings these comparisons into one clinical PBMC evaluation
setting: 361{,}792 paired cells from 40 participants, 293 APT measurements,
and a fixed five-lineage/27-subtype ontology. Five patient-disjoint folds
and subject-balanced pooled Macro-F1 define a common new-patient target.
Classical and neural APT-only models establish reference performance;
Soft-cascade is one hierarchy-aware reference rather than a separate
algorithmic contribution. Our contributions are:
\begin{itemize}
    \item \textbf{An APT cell-typing evaluation resource.}
    A fixed input panel, operational hierarchy, patient folds, explicit scoring
    units and reference predictions support comparisons on the same target.
    Complete data access and portable reproduction remain release requirements.
    \item \textbf{Within-lineage discrimination beyond a candidate-set prior.}
    The full-budget MLP reaches oracle fine F1 of 0.385 versus 0.180 for a
    conditional-proportion baseline. Matched-budget APT/RNA results reveal
    heterogeneous recoverability across lineages, rather than a uniform
    assay-resolution deficit inferred from coarse/fine scores.
    \item \textbf{An RNA-reference increment with a bounded interpretation.}
    Paired APT improves MLP fine F1 by \pairedMLPGainTwoK{} and
    \pairedMLPGainFiveK{} at two RNA widths.
    Its overall patient-by-lineage-shuffle contrast remains unresolved:
    improvement over these RNA references does not yet establish finer-grained
    pairing specificity.
\end{itemize}

The resulting findings concern recoverability under a specified cohort,
taxonomy and training recipe. RNA-derived annotations are not independent
biological truth, and an increment over a restricted RNA representation is
not evidence that RNA intrinsically lacks that information. Patient-disjoint
testing supplies the evaluation foundation, not evidence of blood-based
diagnostic validity.
